% 国赛要求，页面布局上下左右都要留出至少 2.5cm
\usepackage[margin=2.5cm]{geometry}

% 编译检查，必须使用 XeLaTeX 编译
\usepackage{iftex}
\ifXeTeX
    \typeout{=> Compiler Check: XeLaTeX detected. Proceeding.}
\else
    \errmessage{COMPILER ERROR: This document MUST be compiled with XeLaTeX.}
\fi

% 中文支持
\usepackage[fontset=none]{ctex}
\usepackage{xeCJK}
\setcounter{secnumdepth}{5} % 为 \subsubsection 增加编号

% 字体设置
\usepackage{anyfontsize}
\usepackage{fontspec}
\setmainfont{Times New Roman} % 英文使用 Times New Roman
\setmonofont{JetBrains Mono}
\setCJKmonofont{KaiTi}

% 使用基础宋体
\setCJKmainfont{SimSun}[
    % BoldFont = SimHei, % 中文粗体使用黑体
    BoldFont = {思源宋体 CN SemiBold},
    ItalicFont = KaiTi, % 中文斜体使用楷体
    BoldItalicFont = KaiTi
]

% 该定义用于附录仿宋字体
\newCJKfontfamily{\fangsong}{FangSong}

% 章节标题
\usepackage{titlesec}
\renewcommand{\thesection}{\chinese{section}}
\titleformat{\section}
{\Large\bfseries\centering}
{\thesection、}
{0em}
{}

\renewcommand{\thesubsection}{\arabic{section}.\arabic{subsection}}
\titleformat{\subsection}
{\large\bfseries}
{\thesubsection}
{1em}
{}

\renewcommand{\thesubsubsection}{\thesubsection.\arabic{subsubsection}}
\titleformat{\subsubsection}
{\normalsize\bfseries}
{\thesubsubsection}
{1em}
{}

% 标题间距调整
\titlespacing{\section}{0pt}{0.2\baselineskip}{0.2\baselineskip}
\titlespacing{\subsection}{0pt}{0.2\baselineskip}{0.2\baselineskip}
\titlespacing{\subsubsection}{0pt}{0.2\baselineskip}{0.2\baselineskip}

% Drawing 绘图板块包
\usepackage{pgfplots}
\usepackage{tikz}
\usepackage{tikz-cd}
\usepackage{tikz-qtree}
\usepackage{circuitikz}
\usepackage{tikz-3dplot}
\usepackage[svgnames]{xcolor}
\pgfplotsset{compat=1.18}
\usetikzlibrary{shapes.geometric, arrows.meta, positioning, calc, angles, quotes, intersections, tikzmark, backgrounds, shadows, fit}

% 代码块设置
\usepackage{mdframed}
\usepackage{fancyvrb}
\usepackage{listings}
\usepackage{algorithm,algorithmic}

\definecolor{lstcolorComment}{rgb}{0,0.6,0}
\definecolor{lstcolorNumber}{HTML}{E6E6E6}
\definecolor{lstcolorString}{rgb}{1,0.5,0}

\lstset{
    backgroundcolor=\color{white},           % 背景色
    basicstyle=\footnotesize\ttfamily,       % 等宽字体和大小
    breakatwhitespace=false,                 % 是否只在空白处断行
    breaklines=true,                         % 自动断行
    keepspaces=true,                         % 保留代码中的空格，对齐代码
    tabsize=2,                               % Tab 长度为 2 个空格
    frame=single,                            % 单线边框
    rulecolor=\color{black},                 % 边框颜色
    captionpos=bl,                           % 标题位置 (bottom-left)
    numbers=left,                            % 行号在左侧
    numbersep=5pt,                           % 行号与代码的距离
    numberstyle=\tiny\color{lstcolorNumber}, % 行号样式
    stepnumber=1,                            % 行号步长
    commentstyle=\color{lstcolorComment},    % 注释样式
    keywordstyle=\color{blue},               % 关键字样式
    stringstyle=\color{lstcolorString},      % 字符串样式
    showspaces=false,                        % 不显示空格
    showstringspaces=false,                  % 不显示字符串中的空格
    showtabs=false,                          % 不显示 Tab
    escapeinside={\%*}{*},                   % 允许在代码中插入 LaTeX 命令
    extendedchars=true                       % 支持非 ASCII 字符 (对旧编码有效)
}

% Matlab 代码块
\usepackage{matlab-prettifier}
\lstdefinestyle{matlab}{
    style=Matlab-Pyglike,         % 基础样式，提供精确的 MATLAB 关键字高亮
    mlshowsectionrules=true,      % 特色功能：高亮 %% 分节符
    backgroundcolor=\color{white},
    basicstyle=\footnotesize\ttfamily,
    numberstyle=\scriptsize\color{mygray}, % 统一使用 mygray 颜色
    frame=single,
    rulecolor=\color{black},
    captionpos=bl,
    breaklines=true,
    keepspaces=true,
    numbers=left,
    numbersep=8pt,
    showstringspaces=false,
    tabsize=2,
    keywordstyle=\color{deepblue}\bfseries, % 关键字：深蓝、粗体
    commentstyle=\itshape\color{black!35!white}, % 注释：斜体、浅灰
    stringstyle=\color{orange},          % 字符串：橙色
    identifierstyle=\color{myidentifier},  % 标识符(变量/函数名)：深青色
}

% Python 代码块
\usepackage{color}
\definecolor{mygray}{rgb}{0.5,0.5,0.5}
\definecolor{brightgreen}{RGB}{0, 135, 0}
\definecolor{deepblue}{rgb}{0,0,0.5}
\definecolor{deepred}{rgb}{0.6,0,0}
\definecolor{lightblue}{RGB}{91, 192, 255}
\definecolor{orange}{rgb}{1,0.5,0}
\definecolor{myidentifier}{rgb}{0.2,0.4,0.6} % 为标识符高亮深青色
\lstdefinestyle{python}{
    language=Python,
    backgroundcolor=\color{white},          % 背景色
    basicstyle=\footnotesize\ttfamily,      % 代码字体：小号、等宽
    numberstyle=\scriptsize\color{mygray},  % 行号字体：更小号、灰色
    identifierstyle=\color{myidentifier},   % 标识符字体：深蓝
    frame=single,                           % 单线边框
    rulecolor=\color{lightblue},               % 边框颜色
    captionpos=bl,                           % 标题位置：左下角
    keywordstyle=\color{orange}\bfseries, % 关键字：深蓝、粗体
    commentstyle=\itshape\color{black!35!white}, % 设置注释样式，斜体，浅灰
    stringstyle=\color{brightgreen},             % 字符串：橙色
    morekeywords={self},                    % 额外关键字
    emph={MyClass,__init__, GCN, GAT, self},      % 强调词列表
    emphstyle=\color{deepred}\bfseries,     % 强调词样式：深红、粗体
    breaklines=true,                        % 自动断行
    keepspaces=true,                        % 保留空格，用于对齐
    numbers=left,                           % 行号在左
    numbersep=8pt,                          % 行号与代码间距
    showstringspaces=false,                 % 不显示字符串中的空格
    tabsize=2                               % Tab 长度
}

% Darcula Theme
\definecolor{Darcula-bg}{RGB}{43, 43, 43}        % 背景色 (深灰)
\definecolor{Darcula-fg}{RGB}{169, 183, 198}      % 前景色/普通文本 (亮灰)
\definecolor{Darcula-comment}{RGB}{128, 128, 128}   % 注释 (灰色)
\definecolor{Darcula-keyword}{RGB}{204, 120, 50}    % 关键字 (橙色)
\definecolor{Darcula-function}{RGB}{255, 203, 107}  % 函数名 (黄色)
\definecolor{Darcula-string}{RGB}{106, 135, 89}     % 字符串 (绿色)
\definecolor{Darcula-variable}{RGB}{152, 118, 170} % 变量与参数 (淡紫色)
\definecolor{Darcula-number}{RGB}{104, 151, 187}   % 数字 (蓝色)

\lstdefinestyle{darcula}{
    language=Python,
    backgroundcolor=\color{white},
    basicstyle=\small\ttfamily\color{Darcula-fg},
    frame=single,
    rulecolor=\color{Darcula-keyword},
    captionpos=b,
    keywordstyle=\color{Darcula-keyword},
    stringstyle=\color{Darcula-string},
    commentstyle=\itshape\color{Darcula-comment},
    numberstyle=\tiny\color{Darcula-number}, % 行号也用数字的蓝色
    identifierstyle=\color{Darcula-variable},
    emph={read_csv, merge, rename},
    emphstyle=\color{Darcula-function},
    morekeywords=[2]{True, False, None},
    keywordstyle=[2]\color{Darcula-keyword},
    breaklines=true,
    keepspaces=true,
    numbers=left,
    numbersep=10pt,
    showstringspaces=false,
    tabsize=4,
}

% TColorbox
\usepackage[many]{tcolorbox}
\usepackage{fontawesome5}

\definecolor{warningRed}{RGB}{217, 83, 79}
\definecolor{warningRedBack}{RGB}{242, 222, 222}

\newtcolorbox{warning}{
    enhanced,
    shadow={0.8mm}{-0.8mm}{0mm}{black!20!white},
    breakable, % 允许内容跨页
    colback=warningRedBack, % 背景颜色
    colframe=warningRed,    % 边框颜色
    boxrule=0.5pt,          % 边框线宽
    leftrule=2pt,           % 左侧竖线的宽度 (比边框粗，更明显)
    top=2pt,
    bottom=1.8pt,
    right=5em,
    fontupper=\fontsize{9.5pt}{13pt}\selectfont\kaishu,
    fonttitle=\small\bfseries,
    title={
            \faExclamationTriangle[solid]
            \hspace{1.5pt}
            \textbf{Warning}
        },
    coltitle=white,
    attach boxed title to top right={xshift=0pt, yshift=-18.3pt},
    boxed title style={
            colback=warningRed,
            boxrule=0.3pt,
        }
}

\definecolor{tipBlue}{RGB}{91, 192, 222}
\definecolor{tipBlueBack}{RGB}{217, 237, 247}

\newfontfamily{\kaishu}{KaiTi}
\newtcolorbox{tips}{
    enhanced,
    shadow={0.8mm}{-0.8mm}{0mm}{black!20!white},
    breakable,
    colback=tipBlueBack,    % 背景颜色
    colframe=tipBlue,       % 边框颜色
    boxrule=0.5pt,          % 边框线宽
    leftrule=2pt,           % 左侧竖线的宽度
    top=2pt,
    bottom=1.8pt,
    right=5em,
    fontupper=\fontsize{9.5pt}{13pt}\selectfont\kaishu, % 内容字体: {大小}{衬线}
    fonttitle=\small\bfseries,
    title={
            \faLightbulb[solid]
            \hspace{1.5pt}
            \textbf{Tips}
        },
    coltitle=white,
    attach boxed title to top right={xshift=0pt, yshift=-18.3pt},
    % attach boxed title to top left={xshift=10pt, yshift=-10pt},
    boxed title style={
            colback=tipBlue,
            boxrule=0.3pt,
        }
}

% 行距、缩进
\usepackage{setspace}
\usepackage{indentfirst,comment}
\renewcommand{\baselinestretch}{1.3}
\setlength\parindent{2\ccwd} % 中文两字符缩进

% 颜色支持
\usepackage[table]{xcolor}
\definecolor{structurecolor}{RGB}{60,113,183}
\definecolor{main}{RGB}{0,166,82}
\definecolor{second}{RGB}{255,134,24}
\definecolor{third}{RGB}{0,174,247}
\definecolor{winered}{rgb}{0.5,0,0}
\definecolor{bule}{RGB}{18,29,57}
\colorlet{coverlinecolor}{second}

% Math:
\usepackage{array}
\usepackage{sansmath}
\usepackage{amsmath,mathrsfs,amsfonts,amssymb}
\usepackage{esint}
\usepackage{cases}
\usepackage{bm} % 粗体数学符号
\let\Bbbk\relax % 兼容性处理

% 特殊符号
\usepackage{pifont,manfnt,bbding}
\usepackage{adforn}
\usepackage{csquotes} % 引号

% 脚注
\usepackage[flushmargin,stable]{footmisc}
\setlength{\footnotesep}{12pt}

% 图表
\usepackage{tabularx}
\usepackage{booktabs}
\usepackage{multicol,multirow}
\usepackage{makecell}
\usepackage{graphicx}
\usepackage{siunitx}    % 提供 S 列类型，用于对齐数字
\usepackage[hypcap=false, labelfont={bf,color=structurecolor}]{caption}
\usepackage{float}
\usepackage{subcaption}
\captionsetup[table]{skip=0.6\baselineskip, aboveskip=0.6\baselineskip}
\captionsetup[figure]{skip=0.6\baselineskip, aboveskip=0.6\baselineskip}
\captionsetup{
    font=bf,
    labelfont=bf,
    labelsep=quad,
    justification=centering,
    singlelinecheck=true,
    skip=1em
}
\newcommand\figref[1]{\textbf{图}~\ref{#1}}
\newcommand\tabref[1]{\textbf{表}~\ref{#1}}

% 列表，设置了 Itemize 和 Enumerate 不缩进
\usepackage{enumitem}

\setlist[itemize]{
    parsep=\parskip,
    topsep=6pt,
    itemsep=1pt,
    leftmargin=*
}

\setlist[enumerate,1]{
    parsep=\parskip,
    topsep=6pt,
    label=\arabic*., % 标签格式为 1., 2., ...
    leftmargin=* % 标签对齐
}

\setlist[enumerate,2]{
    label=(\alph*),  % 标签格式为 (a), (b), ...
    leftmargin=*,
    topsep=1pt,
    itemsep=0pt
}

% 参考文献
\usepackage[
    backend=biber,          % 指定后端为 biber
    style=gb7714-2015,      % 使用国标 GBT7714-2015 样式
    citestyle=numeric-comp  % 引用为数字压缩格式，如 [1-3]
]{biblatex}
\addbibresource[location=local]{references.bib} % 指定你的 .bib 文件

% \usepackage[numbers,sort&compress]{gbt7714} % 存在多篇引文时自动压缩引号
% \bibliographystyle{gbt7714-numerical}
% \setcitestyle{super=false}
% \renewcommand{\citenumfont}[1]{\textbf{#1}}

% 附录
\usepackage[header]{appendix}
\usepackage{titlesec}

% 只在第一次调用 \addappendix 时执行
\newcommand{\addappendix}[1]{
    \appendix

    \titleformat{\section}
    {\large\bfseries\centering}
    {附录~\thesection}
    {1em}
    {}

    \titlespacing*{\section}{0pt}{2em}{1em}

    % 将 \addappendix 命令重定义为一个简单的 \section 命令
    % 后续再调用 \addappendix 时，就不会重复执行上面的设置了
    % 这里用 ##1 是因为在 \newcommand 的定义中)
    \renewcommand{\addappendix}[1]{\section{##1}}
    \refstepcounter{section}
    \section*{附录~\thesection\hspace{1em}#1}
}
% 此处让 \addappendix 命令暂时处于不进入目录的状态

% 超链接
\usepackage{hyperref}
\hypersetup{
    breaklinks,
    unicode,
    linktoc=all,
    bookmarksnumbered=true,
    bookmarksopen=true,
    colorlinks=true,
    linkcolor=black,
    citecolor=black,
    urlcolor=blue,
    plainpages=false,
    pdfstartview=FitH,
    pdfborder={0 0 0},
    linktocpage
}

\newcommand{\h}{\hspace}
\renewcommand{\v}{\vspace}

\AtBeginDocument{
    \setlength{\abovedisplayskip}{0.45em}
    \setlength{\belowdisplayskip}{0.6em}
}